\documentclass{article}


 \usepackage[final, nonatbib]{neurips_2025}

\usepackage[utf8]{inputenc} % allow utf-8 input
\usepackage[T1]{fontenc}    % use 8-bit T1 fonts
\usepackage{hyperref}       % hyperlinks
\usepackage{url}            % simple URL typesetting
\usepackage{booktabs}       % professional-quality tables
\usepackage{amsfonts}       % blackboard math symbols
\usepackage{nicefrac}       % compact symbols for 1/2, etc.
\usepackage{microtype}      % microtypography
\usepackage{xcolor}         % colors
\usepackage{amsmath}        % equations
\usepackage{graphicx}       % figures
\usepackage{algorithm}      % algorithm float
\usepackage{algpseudocode}  % pseudocode




\title{Structure-Aware Element-level Multimodal Retrieval\\in Academic Documents}


\author{%
    Jaehyeon Lee \\
    20262670\\
    \And
    Seoyeon Lee \\
    20262097\\
    \And
    Suhyeon Jun \\
    20262717\\
    \And
    Minjun Kim \\
    20262242
}


\begin{document}


\maketitle

\begin{abstract}
We propose a \emph{structure-aware} framework for element-level
multimodal retrieval on academic papers, in which the document's
typed element graph serves as a single source of truth. The same
edge-weight schema (i)~defines a graded relevance distribution for a
listwise contrastive loss (\textbf{GRCL}), strictly generalizing binary
InfoNCE, and (ii)~constructs a transition matrix for plug-in
inference-time score propagation. On SPIQA test-A, GRCL improves R@10
by $+2.7$pp and MRR by $+5.4$pp ($p\!=\!0.042$). Most strikingly,
applying the same propagation to the external GME-Qwen2-VL-2B MLLM
raises its R@10 from $58.3$ to $\mathbf{84.4}$ ($+26$pp) and MRR from
$33.1$ to $\mathbf{63.3}$ ($+30$pp) \emph{without any fine-tuning}.
\end{abstract}


\section{Introduction}

Retrieval on academic papers couples two challenges. First, current
multimodal encoders perform unevenly when a text query must surface a
visual element: compact dual encoders such as CLIP and
SigLIP~\cite{siglip2} score poorly on academic figure and table retrieval,
while larger retrieval-tuned MLLMs (e.g., GME-Qwen2-VL-2B~\cite{gme})
reach usable accuracy---suggesting that text-image alignment alone is a
weak signal for this domain. Second, even when the retrieval target is
a single element, it rarely stands alone: a figure is bound to its
caption, the citing paragraphs, and sometimes an appendix reference
through typed relations like \texttt{caption\_of} and \texttt{refer\_to},
so surfacing it benefits from these structural relations even when only
the focal element is annotated as ground truth. Page-level visual
RAG~\cite{colpali,visrag} returns coarse units that cannot pinpoint the
element evidence and tends to miss cross-page references; element-level
retrieval without a graph prior ranks each candidate independently and
ignores the typed relations.

We propose a \emph{structure-aware} retrieval framework in which the
document's typed element graph is a single source of truth, with one
edge-weight schema serving two distinct roles. At training time, a 2-hop max-product path kernel turns the graph
into a graded relevance distribution that supervises a listwise
contrastive loss---the Graph-Relevance Contrastive Loss (GRCL)---which
strictly generalizes binary InfoNCE and recovers it when the targets
degenerate to $\{0,1\}$. At inference time, the same edge weights form a
transition matrix used by weighted diffusion or Personalized PageRank to
smooth any base retriever's scores, with no retraining required.
Empirically, applying this propagation to the external retrieval-tuned
MLLM GME-Qwen2-VL-2B~\cite{gme} raises its SPIQA R@10 from $58.3$ to
$84.4$ ($+26$pp) and almost doubles its MRR ($33.1\!\to\!63.3$), evidence
that the graph functions as an encoder-agnostic universal prior.


\section{Method}

\paragraph{Element graph schema.}
Each document is represented as a graph $G_d\!=\!(V_d,E_d)$ with six
node types (\texttt{text}, \texttt{figure}, \texttt{table},
\texttt{caption}, \texttt{equation}, \texttt{section\_header}) and five
typed edges extracted directly from arXiv LaTeX source ($\backslash$caption
proximity, $\backslash$ref macros, $\backslash$section nesting, and
sequential reading order). The framework's \emph{single source of truth}
is one edge-weight schema:
$\text{BASE}\!=\!\{\texttt{caption\_of}\!:\!1.0,\,
\texttt{refer\_to}\!:\!0.8,\,\texttt{contains}\!:\!0.5,\,
\texttt{reading\_next}\!:\!0.3,\,\texttt{section\_next}\!:\!0.2\}$.
Intuitively, the weights track binding strength: a caption physically
attached to its figure is tightest ($1.0$), an explicit $\backslash$ref
reference is slightly weaker ($0.8$), and section or reading adjacency
weaker still. Two
multiplicative modifiers refine these weights based on the target node:
a section-role modifier (appendix$:\!1.3$, references$:\!0.5$)
emphasizes appendix evidence and de-emphasizes citation lists, and a
cross-modal modifier ($1.1$ for \texttt{refer\_to} into a figure/table)
slightly boosts paragraph-to-figure references over
paragraph-to-paragraph ones---an explicit handle for the
text$\leftrightarrow$visual bridges that make the corpus multimodal.

\paragraph{Encoder.}
We use SigLIPv2-base-patch16-224~\cite{siglip2}, a dual-tower multimodal
encoder with separate text and vision branches, with LoRA rank~8 on the
attention projections (\texttt{q/k/v/out\_proj}) of both towers; the
pretrained backbone is frozen, so only $\sim$491k LoRA + projection-head
parameters are updated ($0.32\%$ of the base). Token-level outputs are projected to
$D_\text{proj}\!=\!128$ and L2-normalized. The retrieval score uses ColBERT-style MaxSim~\cite{colbert},
$\mathrm{score}(q,e)\!=\!\sum_i\max_j\langle q_i,e_j\rangle$.

\paragraph{Graph-Relevance Contrastive Loss (GRCL).}
The kernel between an anchor $q$ and a candidate $e$ is the max-product
weight along the best path of at most two hops, with a per-hop decay
$\gamma\!=\!0.5$:
\begin{equation}
g(q,e)\;=\;\max_{\pi:q\to e,\,|\pi|\le 2}\ \prod_{(u,v)\in\pi} w_{uv}\,\gamma^{|\pi|-1},
\label{eq:kernel}
\end{equation}
where each $w_{uv}$ follows the BASE schema with target-modifier
corrections.
Intuitively, $g(q,e)$ is close to one when $q$ and $e$ are bound by a
strong edge (e.g., a figure and its caption), shrinks on two-hop paths,
and is zero when no short path exists.

For a candidate pool $\{e_1,\ldots,e_K\}$ of size $K\!=\!12$, let
$r_i\!=\!g(q,e_i)$ and $s_i\!=\!\mathrm{score}(q,e_i)$. The loss is a
soft-target listwise cross-entropy,
\begin{equation}
\mathcal{L}_\text{GRCL}(q)\;=\;-\sum_{i=1}^{K}\frac{r_i}{\sum_j r_j}\,\log\frac{\exp(s_i/\tau)}{\sum_j\exp(s_j/\tau)}.
\label{eq:grcl}
\end{equation}
The retriever is therefore asked to match a graded distribution rather
than to pick one positive among uniform negatives---when the targets
degenerate to $\{0,1\}$ this recovers binary InfoNCE, so GRCL is a
strict graded relaxation.
Each anchor is sampled $1/3$ each from \texttt{caption\_of} pairs
(visual anchor), \texttt{refer\_to} pairs (body anchor), and
natural-language QA queries, so a single retriever learns to handle the
three query distributions seen at evaluation. The pool mixes graph
1-hop positives, same-document negatives (hard context that shares
section but lies outside the relevance subgraph), and cross-document
negatives (drawn from other papers in the batch to prevent
within-document shortcuts). Algorithm~\ref{alg:grcl} shows the full
training step.

\begin{algorithm}[t]
\caption{GRCL training step (one anchor)}
\label{alg:grcl}
\begin{algorithmic}[1]
\Require document $d$ with graph $G_d$, encoder $f_\theta$, temperature $\tau$, decay $\gamma$
\State sample anchor $q$ from $\{\texttt{caption\_of},\texttt{refer\_to},\texttt{nl\_qa}\}$ with prob.\ $1/3$ each
\State build pool $C\!=\!\{e_1,\ldots,e_{12}\}$: 1-hop positives $\cup$ same-doc negatives $\cup$ cross-doc negatives
\For{$i = 1,\ldots,K$}
  \State $r_i \gets g(q, e_i)$  \Comment{eq.~(\ref{eq:kernel}); 2-hop max-product BFS on $G_d$}
  \State $s_i \gets \mathrm{MaxSim}(f_\theta(q),\, f_\theta(e_i))$  \Comment{ColBERT-style late interaction}
\EndFor
\State $\mathcal{L} \gets -\sum_i (r_i / \textstyle\sum_j r_j)\log\!\bigl[\,\mathrm{softmax}(s/\tau)\bigr]_i$  \Comment{eq.~(\ref{eq:grcl})}
\State update $\theta$ (LoRA adapter $+$ projection head) via AdamW
\end{algorithmic}
\end{algorithm}

\paragraph{Inference-time graph propagation.}
The intuition is to use the same graph as a smoothing prior: if a
candidate is connected to other high-scoring candidates by strong typed
edges, its own score should rise. Given any base retriever's scores
$s\!\in\!\mathbb{R}^{|V_d|}$, we truncate to the top-$N\!=\!50$
candidates and build a $50\!\times\!50$ transition matrix $P$ by
row-normalizing the \emph{same} edge-weight schema as the training
kernel. We iteratively re-rank with one of three regimes: \emph{uniform
diffusion} sets all edge weights to one (structure-only ablation);
\emph{weighted diffusion} keeps the full BASE weights; both update as
$s^{(t+1)}\!=\!(1\!-\!\alpha)\,s^{(t)}+\alpha\,P^\top s^{(t)}$ for $T$
steps. \emph{Personalized PageRank} re-injects the original scores as a
teleport, $s^{(t+1)}\!=\!\alpha\,q+(1\!-\!\alpha)\,P^\top s^{(t)}$ with
$q\!=\!s^{(0)}$~\cite{ppr}, iterating until
$\|s^{(t)}\!-\!s^{(t-1)}\|_\infty\!<\!10^{-4}$. One iteration is
$<\!1$\,ms, and all three regimes are eval-only and apply to any base
retriever without parameter updates.


\section{Experiment}

\paragraph{Setup.}
SPIQA~\cite{spiqa} is a figure-grounded QA benchmark over arXiv CS papers.
We train on the \texttt{train\_val} split (25{,}459 papers, $\sim$250k
caption--figure pairs) and evaluate on \texttt{test-A} (118 papers,
999 queries). Graphs are extracted directly from arXiv LaTeX source
($\backslash$section, $\backslash$ref) and the dataset's
\texttt{all\_figures} JSON, so no PDF parsing is needed. Training uses
AdamW for $8{,}000$ steps with $16$ anchors $\times\,12$ candidates per
step; the default setting is $\text{lr}\!=\!3$e-5, $\tau\!=\!0.07$,
$\gamma\!=\!0.5$, and a tuned setting is $\text{lr}\!=\!1$e-4, $\tau\!=\!0.10$.
One row takes $\sim$5~h on a single A100 80GB. Evaluation uses a
per-document candidate pool, reporting R@5, R@10, and MRR; significance is
assessed by paired bootstrap ($1{,}000$ resamples). We compare against
two baselines: the same encoder trained with binary InfoNCE, and the
external retrieval-tuned MLLM GME-Qwen2-VL-2B-Instruct~\cite{gme} used
zero-shot.

\paragraph{Results.}
Table~\ref{tab:main} summarizes our findings; specific rows are
referenced below.
GRCL outperforms binary InfoNCE on in-domain SPIQA (rows~1 vs.~2) by
$+2.7$pp R@10 and $+5.4$pp MRR; paired bootstrap on Coverage@10 gives
$\Delta\!=\!+0.0255$, $95\%$ CI $[+0.0015,+0.0511]$, $p\!=\!0.042$. The
relative MRR gain ($+16.5\%$) is roughly $5\times$ the R@10 gain,
indicating that graded supervision sharpens top-rank quality. Restricting
the GRCL kernel to a single edge type sharpens the signal further:
\texttt{refer\_to}-only edges (row~3) yield R@10 $=\mathbf{90.4}$, the
best trained-encoder result in our ablation. A $\backslash$ref macro
is a manual evidence annotation by the paper's author, whereas
\texttt{caption\_of} is a trivial layout co-occurrence; mixing weak
edges dilutes the strong signal and costs $3.3$pp R@10 versus the
multi-edge GRCL kernel. Hyperparameter tuning is a comparable
lever (rows~4--5): raising $\text{lr}$ to $1$e-4 reaches R@10 $88.6$ and
MRR $42.6$ (row~4), and the full HP combo of $\text{lr}\!+\!\tau\!+\!\lambda$
reaches R@10 $88.0$ and MRR $42.2$ (row~5).

The single largest effect is observed not in our trained model but as a
zero-shot plug-in (row~7 vs.~8--12): applying the same edge schema on
top of GME-Qwen2-VL-2B raises R@10 from $58.3$ to $84.4$ ($+26$pp) and
MRR from $33.1$ to $63.3$ ($+30$pp) without any fine-tuning, reaching
$96\%$ of our best trained SigLIP's R@10. The three regimes are
equivalent on GME (rows~8--12)---uniform diffusion, weighted diffusion,
and PPR all converge to R@10 $83$--$85\%$ and MRR $59$--$63\%$---
indicating that propagation utility comes from graph structure itself
rather than from modifier sophistication.

\begin{table}[t]
\centering
\small
\caption{SPIQA test-A retrieval. Bold marks the column best or the
headline finding; $\Delta$ values are versus the in-group baseline (row~1
for trained encoders; row~7 for GME).}
\label{tab:main}
\begin{tabular}{lccccc}
\toprule
Setup & R@5 & R@10 & MRR & $\Delta$R@10 & $\Delta$MRR \\
\midrule
\multicolumn{6}{l}{\emph{Trained dual-encoder: SigLIPv2-base + LoRA, GRCL unless noted}} \\
1.\ InfoNCE baseline & --- & 84.4 & 32.8 & --- & --- \\
2.\ GRCL ($p\!=\!0.042$ vs.\ 1) & --- & 87.1 & 38.2 & $+2.7$ & $+5.4$ \\
3.\ \quad+ \texttt{refer\_to}-only edges & --- & \textbf{90.4} & 41.9 & $\mathbf{+6.0}$ & $+9.1$ \\
4.\ \quad+ $\text{lr}\!=\!1$e-4 & --- & 88.6 & \textbf{42.6} & $+4.2$ & $\mathbf{+9.8}$ \\
5.\ \quad+ best HP combo (lr$+\tau+\lambda$) & --- & 88.0 & 42.2 & $+3.6$ & $+9.4$ \\
6.\ \quad+ CLIP-L/14 backbone swap & --- & 88.1 & 36.2 & $+3.7$ & $+3.4$ \\
\midrule
\multicolumn{6}{l}{\emph{External MLLM: GME-Qwen2-VL-2B-Instruct, zero-shot, no fine-tuning}} \\
7.\ GME, no propagation & 47.1 & 58.3 & 33.1 & --- & --- \\
8.\ \quad+ uniform diffusion ($\alpha\!=\!0.5,\,T\!=\!2$) & --- & 83.9 & 62.1 & $+25.6$ & $+29.0$ \\
9.\ \quad+ weighted diffusion ($\alpha\!=\!0.3,\,T\!=\!2$) & \textbf{78.2} & 84.4 & 63.3 & $+26.1$ & $+30.2$ \\
10.\ \quad+ weighted diffusion ($\alpha\!=\!0.3,\,T\!=\!3$) & --- & \textbf{84.5} & \textbf{63.4} & $\mathbf{+26.2}$ & $\mathbf{+30.3}$ \\
11.\ \quad+ PPR ($\alpha\!=\!0.30$ teleport) & --- & 83.3 & 59.1 & $+25.0$ & $+26.0$ \\
12.\ \quad+ PPR ($\alpha\!=\!0.50$) & --- & 83.2 & 58.6 & $+24.9$ & $+25.5$ \\
\bottomrule
\end{tabular}
\end{table}


\section{Conclusion}

A typed document element graph can simultaneously serve as graded
supervision for retriever training and as an inference-time prior for any
base retriever, driven by a single edge-weight schema. On SPIQA, GRCL
gives a statistically significant gain over binary InfoNCE,
\texttt{refer\_to}-only edges yield the best trained R@10 ($90.4$), and
plug-in propagation atop the external GME-Qwen2-VL-2B MLLM produces our
largest single effect ($+26$pp R@10, $+30$pp MRR).


\newpage


\bibliographystyle{splncs04}
\bibliography{reference}

\end{document}
